% ===============================================================================
% AB-04d: Wiederholung - Mi colegio
% Spanisch Klasse 7 - Sequenz 4: Mi colegio
% ===============================================================================
% Abgeleitet von: LE-04d.tex (Score: 9.1/10)
% Punkte: 20 (4+4+6+6)
% ===============================================================================

\documentclass[11pt, a4paper]{article}

% ===============================================================================
% SW-MODUS TOGGLE
% ===============================================================================
\newif\ifswmodus
\swmodusfalse

% ===============================================================================
% PAKETE
% ===============================================================================

\usepackage[utf8]{inputenc}
\usepackage[T1]{fontenc}
\usepackage[ngerman]{babel}
\usepackage{helvet}
\renewcommand{\familydefault}{\sfdefault}

\usepackage[margin=2cm]{geometry}
\usepackage{enumitem}
\usepackage{tabularx}
\usepackage{booktabs}
\usepackage{multirow}

\usepackage{xcolor}
\usepackage{framed}
\usepackage{tcolorbox}
\tcbuselibrary{skins}

\usepackage{fancyhdr}
\usepackage{lastpage}
\usepackage{needspace}

% ===============================================================================
% FARBEN
% ===============================================================================

\definecolor{spanisch-color}{HTML}{c41e3a}
\definecolor{grau-color}{HTML}{888888}
\definecolor{hellgrau-color}{HTML}{f5f5f5}
\definecolor{orange-color}{HTML}{b35c00}

\definecolor{sw-dark}{HTML}{000000}
\definecolor{sw-gray}{HTML}{666666}
\definecolor{sw-white}{HTML}{ffffff}

\ifswmodus
    \colorlet{spanisch}{sw-dark}
    \colorlet{grau}{sw-gray}
    \colorlet{hellgrau}{sw-white}
    \colorlet{orange}{sw-dark}
\else
    \colorlet{spanisch}{spanisch-color}
    \colorlet{grau}{grau-color}
    \colorlet{hellgrau}{hellgrau-color}
    \colorlet{orange}{orange-color}
\fi

\newcommand{\fachfarbe}{spanisch}

% ===============================================================================
% VARIABLEN
% ===============================================================================

\newcommand{\thema}{Wiederholung: Mi colegio}
\newcommand{\sequenz}{04}
\newcommand{\block}{d}
\newcommand{\fach}{Spanisch}
\newcommand{\klasse}{7}
\newcommand{\maxpunkte}{20}
\newcommand{\datum}{\today}

% ===============================================================================
% KOPF- UND FUSSZEILE
% ===============================================================================

\pagestyle{fancy}
\fancyhf{}
\renewcommand{\headrulewidth}{0.4pt}
\renewcommand{\footrulewidth}{0.4pt}

\lhead{\fach{} -- Klasse \klasse}
\chead{AB-\sequenz\block}
\rhead{\datum}

\lfoot{}
\cfoot{}
\rfoot{Seite \thepage\ von \pageref{LastPage}}

% ===============================================================================
% BEFEHLE
% ===============================================================================

\newcommand{\punktebox}[1]{\hfill\fbox{\small \_/#1}}
\newcommand{\luecke}[1]{\rule{#1}{0.4pt}}

\newtcolorbox{merkkasten}[1][]{
    colback=hellgrau,
    colframe=\fachfarbe,
    colbacktitle=white,
    coltitle=\fachfarbe,
    boxrule=1pt,
    arc=0pt,
    left=8pt, right=8pt, top=8pt, bottom=8pt,
    fonttitle=\bfseries,
    attach boxed title to top left={yshift=-2mm, xshift=5mm},
    boxed title style={boxrule=0pt, colback=white},
    title=#1
}

\newtcolorbox{vokabelkasten}{
    colback=white,
    colframe=\fachfarbe,
    boxrule=0.5pt,
    arc=0pt,
    left=5pt, right=5pt, top=5pt, bottom=5pt
}

\newtcolorbox{hinweis}{
    colback=white,
    colframe=orange,
    colbacktitle=white,
    coltitle=orange,
    boxrule=1pt,
    arc=0pt,
    left=5pt, right=5pt, top=5pt, bottom=5pt,
    fonttitle=\bfseries,
    attach boxed title to top left={yshift=-2mm, xshift=5mm},
    boxed title style={boxrule=0pt, colback=white},
    title=Hinweis
}

\newenvironment{aufgabe}[1]{%
    \needspace{5\baselineskip}%
    \nopagebreak[4]%
    \vspace{10pt}
    \noindent\textbf{\textcolor{\fachfarbe}{Aufgabe #1}}
    \nopagebreak[4]%
    \vspace{5pt}
    \nopagebreak[4]%
    \begin{list}{}{\leftmargin=0pt}
    \item[]
}{%
    \end{list}
}

\newcommand{\es}[1]{\textcolor{\fachfarbe}{\textbf{#1}}}

% ===============================================================================
% DOKUMENT
% ===============================================================================

\begin{document}

% ===============================================================================
% KOPFBEREICH
% ===============================================================================

\begin{center}
    {\Large\bfseries\textcolor{\fachfarbe}{Arbeitsblatt: \thema}}
\end{center}

\vspace{5pt}

\noindent
\begin{tabularx}{\textwidth}{|X|X|X|}
\hline
\textbf{Name:} \luecke{4cm} &
\textbf{Klasse:} \klasse &
\textbf{Datum:} \luecke{2cm} \\
\hline
\end{tabularx}

\vspace{15pt}

% ===============================================================================
% MERKKASTEN: Schulfächer
% ===============================================================================

\begin{merkkasten}[Merkkasten: Las asignaturas (Schulfächer)]

\begin{tabular}{ll|ll}
\es{las matemáticas} & = \luecke{3cm} & \es{la música} & = \luecke{3cm} \\[0.5em]
\es{el inglés} & = \luecke{3cm} & \es{la historia} & = \luecke{3cm} \\[0.5em]
\es{el español} & = \luecke{3cm} & \es{la educación física} & = \luecke{3cm} \\[0.5em]
\es{el arte} & = \luecke{3cm} & \es{las ciencias naturales} & = \luecke{3cm} \\
\end{tabular}

\end{merkkasten}

\vspace{10pt}

% ===============================================================================
% AUFGABE 1: Schulfächer zuordnen
% ===============================================================================

\begin{aufgabe}{1 -- Schulfächer zuordnen}
Verbinde das spanische Schulfach mit der deutschen Bedeutung. \punktebox{4}
\end{aufgabe}

\begin{center}
\begin{tabular}{rcl}
\es{las matemáticas} & \luecke{2cm} & Sport \\[0.8em]
\es{la música} & \luecke{2cm} & Mathematik \\[0.8em]
\es{la historia} & \luecke{2cm} & Musik \\[0.8em]
\es{la educación física} & \luecke{2cm} & Geschichte \\
\end{tabular}
\end{center}

\vspace{15pt}

% ===============================================================================
% AUFGABE 2: Wochentage + Uhrzeiten
% ===============================================================================

\begin{aufgabe}{2 -- Wochentage und Uhrzeiten}
Ergänze die fehlenden Wörter. \punktebox{4}
\end{aufgabe}

\noindent
\textbf{a) Wochentage:}

\vspace{0.5em}

\noindent
lunes, \luecke{2.5cm}, miércoles, \luecke{2.5cm}, viernes, \luecke{2.5cm}, domingo

\vspace{1em}

\noindent
\textbf{b) Uhrzeiten -- Übersetze:}

\vspace{0.5em}

\noindent
Um 8 Uhr = \luecke{4cm} \hfill Um halb 10 = \luecke{4cm}

\vspace{15pt}

% ===============================================================================
% MERKKASTEN: -ar Verben
% ===============================================================================

\begin{merkkasten}[Merkkasten: Verben auf -ar konjugieren]

\begin{center}
\begin{tabular}{|l|l|l|}
\hline
\textbf{Person} & \textbf{estudiar} & \textbf{Endung} \\
\hline
yo & \luecke{2.5cm} & -o \\
\hline
tú & \luecke{2.5cm} & -as \\
\hline
él/ella & \luecke{2.5cm} & -a \\
\hline
nosotros & \luecke{2.5cm} & -amos \\
\hline
ellos/ellas & \luecke{2.5cm} & -an \\
\hline
\end{tabular}
\end{center}

\end{merkkasten}

\vspace{10pt}

% ===============================================================================
% AUFGABE 3: -ar Verben konjugieren
% ===============================================================================

\begin{aufgabe}{3 -- Konjugiere die Verben}
Setze die richtige Form ein. \punktebox{6}
\end{aufgabe}

\noindent
a) Yo \luecke{2.5cm} (hablar) español. \hfill
b) Tú \luecke{2.5cm} (cantar) bien.

\vspace{0.8em}

\noindent
c) Ella \luecke{2.5cm} (estudiar) matemáticas. \hfill
d) Nosotros \luecke{2.5cm} (escuchar) música.

\vspace{0.8em}

\noindent
e) Ellos \luecke{2.5cm} (terminar) a las dos. \hfill
f) Yo \luecke{2.5cm} (bailar) salsa.

\vspace{15pt}

% ===============================================================================
% AUFGABE 4: Schulalltag beschreiben
% ===============================================================================

\begin{aufgabe}{4 -- Describe tu día en el colegio}
Schreibe 3 Sätze über deinen Schulalltag. Verwende Schulfächer, Wochentage und Uhrzeiten. \punktebox{6}
\end{aufgabe}

\noindent
\textbf{Beispiel:} Los lunes tengo matemáticas a las ocho.

\vspace{0.8em}

\noindent
1. \luecke{12cm}

\vspace{0.8em}

\noindent
2. \luecke{12cm}

\vspace{0.8em}

\noindent
3. \luecke{12cm}

% ===============================================================================
% PUNKTEÜBERSICHT
% ===============================================================================

\vspace{20pt}
\hrule
\vspace{10pt}

\begin{flushright}
\textbf{Punkte:} \luecke{1cm} / \maxpunkte
\end{flushright}

\end{document}
